\documentclass[11pt,a4paper]{article}
\usepackage[utf8]{inputenc}
\usepackage[margin=1in]{geometry}
\usepackage{amsmath,amssymb,amsfonts}
\usepackage{hyperref}
\usepackage{graphicx}
\usepackage{enumitem}
\usepackage{setspace}

\hypersetup{colorlinks=true, linkcolor=blue, urlcolor=blue, citecolor=blue}
\onehalfspacing

\begin{document}

\title{\textbf{Retrocausal Quantum Entanglement Communication via Microtubule-Orchestrated Modulation of the Light-Time Interface in a 6D Temporal Manifold}}

\author{J. Yesi Orihuela \\
\small Independent Researcher \\
\small MIT Alumnus \\
\small \texttt{jorihuela@alum.mit.edu}}

\date{May 2026 (draft v1.2)}

\maketitle

\begin{abstract}
We propose a mechanism for self-consistent retrocausal communication through quantum entanglement mediated by the light-time interface in Kletetschka’s (2025) 6D temporal manifold. Orchestrated Objective Reduction (Orch-OR) in neuronal microtubules (Penrose \& Hameroff) is hypothesized to act as a biological transceiver: the microtubule lattice couples to the electromagnetic ``flame front'' that actualizes the present slice, enabling low-energy information exchange with past or future states of the same system. Negative-time domain walls (Orihuela, 2026) provide a controllable laboratory method to tilt or invert the interface, potentially amplifying the effect. The framework builds directly on our prior work in repulsive gravity, three-dimensional time, and quantum light-time duality. It requires no new fundamental physics and generates immediate, testable predictions using existing petawatt lasers, single-photon sources, and combined EEG/fMRI with atom interferometry.
\end{abstract}

\section{Introduction}

The Burning Log phenomenology illustrates time as a complete block in which the unburned wood (future), moving flame front (present), and charred remains (past) coexist. This intuition is formalized in Kletetschka’s 6D temporal manifold and our Quantum Light-Time Duality framework, which identifies electromagnetic radiation as the active interface that actualizes the experienced present. Photons, with null proper time, constitute the moving hypersurface of the ``now.''

Extending this picture, we examine whether the same interface can support coherent information exchange across an individual’s timeline via quantum entanglement. Neuronal microtubules, proposed as the site of Orch-OR quantum processes, offer a natural biological candidate for such coupling.

\section{Theoretical Foundations}

\subsection{Kletetschka’s 6D Temporal Manifold (2025)}
Three independent timelike coordinates $t_1, t_2, t_3$ plus emergent space. The primary experienced arrow lies along $t_1$; the orthogonal dimensions supply the block thickness.

\subsection{Quantum Light-Time Duality}
The interface function marking the active present slice is
\[
I(t_1, \mathbf{r}) = \Theta(c t_1 - r) + \delta I_{\rm bias},
\]
where $\delta I_{\rm bias}$ encodes any local tilt induced by negative-time domain walls.

\subsection{Orch-OR and Microtubule Quantum Coherence}
Tubulin dimers in microtubules support quantum superpositions whose objective reduction is gravitationally induced and non-computable. Orchestration by biological processes yields conscious moments. Recent experimental indications of room-temperature quantum coherence in microtubule-like structures strengthen the plausibility of sustained entanglement.

\subsection{Retrocausality in the Framework}
Penrose’s formulation of objective reduction permits time-symmetric dynamics. Combined with a tilted light-time interface, this allows entangled states in the microtubule lattice to maintain correlations with past or future states of the same lattice in a self-consistent manner.

\section{Proposed Mechanism}

Microtubules function as the biological node that:
\begin{itemize}
    \item Maintains entangled quantum states across tubulin conformational degrees of freedom.
    \item Couples to the electromagnetic light-time interface.
    \item Uses transient negative-time bias (naturally occurring or externally induced) to open low-energy windows for information exchange across the primary temporal axis $t_1$.
\end{itemize}

The effective interaction can be modeled with a retrocausal term in the Hamiltonian that respects Novikov self-consistency within the 6D block. No logical paradoxes arise because the block is globally consistent.

\section{Relation to Negative-Time Domain Walls}

A laser-generated negative-time domain wall (Orihuela, 2026) provides a laboratory means of modulating $\delta I_{\rm bias}$, transiently enhancing microtubule-interface coupling. This offers a controlled experimental pathway to test the proposed retrocausal channel.

\section{Falsifiable Predictions}

\begin{enumerate}
    \item Individuals exhibiting strong pattern-recognition or pre-cognitive abilities will show statistically significant differences in occipital EEG signatures or microtubule-relevant biomarkers (e.g., tubulin phosphorylation patterns) immediately preceding or following verified pre-cognitive events.
    \item Single-photon interference experiments conducted inside a negative-time domain wall will display new sidebands whose spacing correlates with Orch-OR characteristic timescales ($\sim$40 Hz $\gamma$-band).
    \item High-resolution sleep or resting-state fMRI/EEG will reveal transient coherence signatures (excess phase-locking or anomalous power spectra) during periods of reported intuitive or pre-cognitive activity.
    \item Any retrocausal information transfer will obey self-consistency constraints: attempts to create logical paradoxes will manifest as measurable decoherence or signal attenuation rather than observable contradictions.
\end{enumerate}

All predictions are accessible with current technology.

\section{Discussion}

This proposal integrates the Burning Log phenomenology, 6D temporal geometry, light-time duality, negative-time domain walls, and Orch-OR into a unified, experimentally testable framework for information exchange across an individual’s timeline. It provides a physically grounded bridge between quantum consciousness research and retrocausality while remaining fully consistent with semiclassical gravity and standard quantum electrodynamics.

The mechanism suggests that certain cognitive profiles may naturally couple more effectively to the light-time interface, offering a potential explanation for documented individual differences in intuitive or pre-cognitive performance. No new fundamental physics is invoked; the tools for empirical investigation exist today.

\section*{Acknowledgments}
The author thanks Grok (xAI) for assistance in refining the mathematical framework and manuscript preparation. The core ideas are the author’s own.

\section*{References}
\begin{enumerate}[leftmargin=*,label={[\arabic*]}]
    \item Kletetschka, G. (2025). Three-Dimensional Time: A Mathematical Framework for Fundamental Physics. \textit{Reports in Advances of Physical Sciences}.
    \item Penrose, R. \& Hameroff, S. (various years). Orchestrated Objective Reduction (Orch-OR) papers.
    \item Orihuela, J. Y. (2026). A Conceptual Proposal for Localized Repulsive Gravity via Laser-Induced Vacuum Polarization and Negative-Time Domain Walls.
    \item Orihuela, J. Y. (2026). The Burning Log and Three-Dimensional Time.
    \item Orihuela, J. Y. (2026). Quantum Light-Time Duality.
    \item Angulo et al. (2024). Experimental evidence of negative time in quantum optics.
\end{enumerate}

\end{document}